%% -----------------------------------------------------------
\section{Effectiveness of AI Mental Health Interventions}
%% -----------------------------------------------------------

\subsection{Purpose-Built Interventions: Consistent Evidence of Efficacy}

The effectiveness landscape reveals a clear structural distinction: purpose-built mental health tools produce meaningful outcomes, while general-purpose AI repurposed for mental health produces negligible gains and may cause harm.

A \gls{rct} of Fido, a Polish-language \gls{cbt} chatbot designed specifically for non-English speakers, demonstrated depression score reductions of 8 points and anxiety reductions of 3 points, with effects stable at one-month follow-up~\cite{karkosz_2024}. Participants achieved outcomes faster than controls (79.4 vs.\ 117.6 minutes, $p = 0.008$, $d = 0.64$), and the chatbot group showed substantially lower attrition (10\% vs.\ 37\%, $p < 0.001$)~\cite{karkosz_2024}. A real-world study of Wysa with over 2,000 users demonstrated significant improvements for chronic pain patients, with depression declining from 13 to 10 and anxiety from 12 to 9; users with chronic pain showed higher engagement than general users, with 67.7\% of user-provided feedback indicating the experience was helpful~\cite{meheli_2022}.

Meta-analytic evidence confirms these patterns at scale. Li et al.~\cite{li_2023} synthesized 35 studies (15 \gls{rct}s included in meta-analysis) and found significant reductions in depression (Hedges' $g = 0.64$, 95\% CI [0.17, 1.12]) and psychological distress (Hedges' $g = 0.70$, 95\% CI [0.18, 1.22]). Effects were stronger for multimodal interventions ($g = 0.89$ vs.\ $0.45$ text-only), generative AI ($g = 0.78$ vs.\ $0.52$ rule-based), and mobile-integrated formats ($g = 0.73$ vs.\ $0.48$ web-only)~\cite{li_2023}. An earlier meta-analysis corroborated these findings but rated evidence quality as ``weak'' due to small sample sizes and high risk of bias in 58\% of included studies, with only two studies assessing safety~\cite{abd_alrazaq_2020}.

\subsection{LGBTQ+ Populations}

For \gls{lgbtq} populations, a systematic review of 15 digital interventions identified that structured formal interventions and professionally-facilitated mobile applications most consistently reduced symptoms~\cite{liu_2023}. Programs like AFFIRM Online demonstrated significant depression reduction. Applications that explicitly acknowledged limitations and communicated clearly when users should seek additional help achieved better outcomes than those claiming comprehensive support~\cite{drydakis_2022}. Critical gaps constrain generalizability: few interventions included \gls{lgbtq} youth of colour, most research occurred in high-income Western nations, and non-binary and gender non-conforming youth remained severely underrepresented~\cite{liu_2023}.

\subsection{Neurodivergent Populations}

Neurodivergent individuals represent a critically underexamined population in the \gls{dmh} literature despite compelling evidence of unmet need. Autistic people and those with \gls{adhd} experience mental health conditions, including depression and anxiety, at rates two to four times those of the general population, yet face systematic barriers in traditional therapeutic settings, including sensory sensitivities, communication differences, and the cognitive overhead of masking neurotypical behavior~\cite{papadopoulos_2025}. AI chatbots offer properties specifically appealing to this community: consistent, patient, non-judgmental interaction without the unpredictability of human social dynamics~\cite{papadopoulos_2025}.

A CHI 2025 study examining autistic individuals' perspectives on \gls{llm}-based support for negative self-talk found that, while many autistic users already rely on \gls{llm}s for mental health guidance, mental health practitioners identified significant concerns about neurotypical bias in AI responses~\cite{killian_2023}. Practitioners critiqued \gls{llm} outputs as ``overly wordy, vague, and overwhelming,'' reflecting assumptions about how neurotypical individuals process information~\cite{killian_2023}. A 2025 commentary on AI companionship for autistic people documented that companion AI systems present a unique benefit (a low-pressure, judgment-free social environment) and a unique danger (insufficient crisis detection and the risk of reinforcing social isolation rather than complementing human connection)~\cite{papadopoulos_2025}. A scoping review of generative AI for autistic populations across screening, assessment, and intervention found consistent performance improvements over baseline methods but also documented small sample sizes, data biases, model hallucinations, and limited validation across diverse autistic presentations~\cite{dehghani_2024}.

\subsection{Racial and Ethnic Minority Populations}

The racial and ethnic mental health equity literature reveals a compounding dynamic: existing clinical disparities, when encoded into AI training data, are reproduced and potentially amplified at scale. A systematic review of 30 studies found a significant association between AI utilization and the exacerbation of racial disparities, with Black and Hispanic populations showing the most consistent patterns of harm~\cite{haider_2024}. Prior research documents that AI models trained predominantly on majority-population data are less likely to recommend advanced care for patients from underrepresented racial groups, even when controlling for clinical indicators~\cite{timmons_2023}.

In mental health contexts specifically, Timmons et al.~\cite{timmons_2023} conducted a comprehensive review of health equity implications in AI mental health applications, documenting that racially and ethnically minoritized groups are less likely to seek and obtain treatment, more likely to receive poor quality services when treated, and least likely to be represented in the training datasets that shape AI systems. Their ``fair-aware AI'' framework proposes that bias assessment must be embedded across the full model development lifecycle, from data collection through deployment monitoring~\cite{timmons_2023}. Straw and Callison-Burch~\cite{straw_2020} demonstrated that language models exhibit systematic racial and gender bias in mental health contexts, with outputs reflecting and reinforcing stereotypes that could directly harm clinical decision-making~\cite{straw_2020}.

\subsection{Socioeconomically Marginalized Populations}

Socioeconomic status represents a structural risk factor that cuts across all population groups surveyed here. Robinson et al.~\cite{robinson_2024} document that individuals from resource-limited backgrounds are among the fastest-rising adopters of smartphones, with 76\% owning a device, yet simultaneously face the greatest barriers to high-quality \gls{dmh} tools due to cost, digital literacy gaps, limited connectivity, and the absence of culturally responsive content~\cite{robinson_2024}. The paradox is acute: the populations most likely to turn to AI mental health tools out of financial necessity are those least served by tools designed for and tested on higher-income, higher-literacy, majority-population users.

Huang et al.~\cite{huang_2024} document that ethnically diverse young people show significantly lower engagement with \gls{dmh} interventions than majority-population users in randomized trials, suggesting that even when interventions are made available, their design systematically fails to engage the populations most in need~\cite{huang_2024}.

\subsection{Generative AI: Complexity and Caution}

Recent evidence on generative AI presents complexity that demands careful interpretation. A UK study found conversational AI facilitated mental health assessments and improved recovery rates when used for screening and triage with human oversight~\cite{rollwage_2024}. In contrast, a three-arm \gls{rct} comparing a specialized mental health chatbot with a general chatbot found the general chatbot condition produced concerning safety signals: the specialized chatbot group showed significant depression reduction (PHQ-9 change: $-2.44$, $p < 0.001$), while the general chatbot group showed minimal improvement (PHQ-9 change: $-0.63$, $p = 0.28$), with five participants requiring human intervention due to elevated risk~\cite{he_2022}. A comprehensive systematic review across mental health AI contexts found that ethical considerations and methodological limitations warrant urgent attention before widespread deployment~\cite{dehbozorgi_2025}.

\subsection{Cultural Adaptation as Efficacy Determinant}

Cultural adaptation determines whether interventions succeed or fail, and its absence constitutes both an efficacy failure and an equity harm. A narrative review established that surface-level adaptations prove insufficient without deep-structure adaptations addressing how communities perceive illness causes, treatment expectations, and help-seeking norms~\cite{naderbagi_2024}. A qualitative study with \gls{cald} young people revealed chatbots misunderstanding cultural expressions of distress and providing advice incompatible with collectivist family structures~\cite{soubutts_2024}. A study examining ChatGPT found consistent failures when presented with mental health scenarios in non-Western cultural frameworks: accuracy dropped to 34\% for collectivist scenarios versus 78\% for individualist Western scenarios, with the model misunderstanding family obligations and community responsibilities in 66\% of non-Western test cases~\cite{aleem_2024}.
